\begin{table}[t]
\centering
\caption{VAT Gap by Country, 2013--2022 (\% of VTTL)}
\label{tab:vatgap}
\small
\begin{tabular}{lcccc}
\hline\hline
Year & Lithuania & Latvia & Estonia & EU Average \\
\hline
2013 & 36.0 & 18.2 & 7.1 & 15.2 \\
2014 & 34.6 & 17.8 & 6.8 & 14.1 \\
2015 & 25.6 & 17.2 & 5.5 & 13.2 \\
2016 & 24.5 & 12.0 & 3.2 & 12.2 \\
2017 & 22.7 & 11.6 & 4.5 & 11.2 $\leftarrow$ \textit{i.SAF} \\
2018 & 24.3 & 11.5 & 2.3 & 11.2 \\
2019 & 21.4 & 8.2 & 3.2 & 10.3 \\
2020 & 12.6 & 6.1 & 0.4 & 9.1 \\
2021 & 5.2 & 8.6 & 5.3 & 5.4 \\
2022 & 1.3 & 6.2 & 1.8 & 5.3 \\
\hline
Pre--2017 mean & 30.2 & 16.3 & 5.7 & 13.7 \\
Post--2016 mean & 14.6 & 8.7 & 2.9 & 8.8 \\
$\Delta$ & $-15.6$ & $-7.6$ & $-2.7$ & $-4.9$ \\
\hline\hline
\end{tabular}
\begin{tablenotes}
\item \textit{Notes:} VAT gap as a percentage of VAT Total Tax Liability (VTTL), from the European Commission/CASE VAT Gap Study Reports (2014--2023 editions). Lithuania's i.SAF mandatory invoice reporting became effective October 1, 2016. The arrow marks the first full post-treatment year.
\end{tablenotes}
\end{table}
